\documentclass[12pt,a4paper]{article}
        \makeatletter
        \def\input@path{{../}}
        \makeatother
        \usepackage{almeria-fichas}

        \FichaMeta{Sesion 02}{¿Que se puede medir y representar en una ciudad?}{Bloque 1 · Mapa y coordenadas}{Distinguir entre ubicacion, medida y dato para organizar correctamente la informacion de Almeria.}{Hechos}{Recordar, Comprender, Aplicar}

        \begin{document}

        \FichaCabecera

        \begin{materialesbox}
        \begin{itemize}
    \item Ficha impresa
    \item Lapiz
    \item Regla pequena
    \item Lapis de tres colores
\end{itemize}
        \end{materialesbox}

        \begin{pasosbox}
        \begin{enumerate}
    \item Lee cada ejemplo de ciudad con calma.
    \item Rodea en azul las ubicaciones, en verde las medidas y en naranja los datos.
    \item Comprueba con un companero una respuesta antes de pasar al siguiente ejercicio.
    \item Escribe siempre la unidad cuando aparezca una medida.
\end{enumerate}
        \end{pasosbox}

        \begin{recordatoriobox}
        \begin{itemize}
    \item Ubicar es decir donde esta algo.
    \item Medir es expresar una magnitud con numero y unidad.
    \item Representar es mostrar la informacion en tabla, plano o grafica.
\end{itemize}
        \end{recordatoriobox}

        \begin{apoyobox}
        \begin{itemize}
    \item Usa una tabla de tres columnas si necesitas ordenar ideas.
    \item Si dudas, pregúntate: ¿hay numero? ¿hay unidad? ¿hay lugar?
    \item Haz una marca al lado de cada dato ya clasificado.
\end{itemize}
        \end{apoyobox}

        \BloomSection{recordar}

\BloomExercise{recordar}{1}{Clasificacion rapida}{Clasifica estas expresiones en \textbf{ubicacion}, \textbf{medida} o \textbf{dato}: \emph{Plaza Vieja}, \emph{250 m}, \emph{18 ºC}, \emph{Rambla}, \emph{3 km}.}{5}

\BloomExercise{recordar}{2}{Vocabulario clave}{Define con una frase sencilla estas palabras: \textbf{ubicacion}, \textbf{medida}, \textbf{dato}.}{5}

\BloomSection{comprender}

\BloomExercise{comprender}{1}{Explica la diferencia}{Explica la diferencia entre estas dos frases: \emph{La Catedral esta en el centro} y \emph{La torre de la Catedral mide 55 m}.}{6}

\BloomExercise{comprender}{2}{Representacion adecuada}{Relaciona cada informacion con una forma de representarla mejor: \emph{lista}, \emph{tabla}, \emph{plano}. Justifica una de tus decisiones.}{6}

\BloomSection{aplicar}

\BloomExercise{aplicar}{1}{Tabla de ciudad}{Completa una tabla de tres columnas con \textbf{dos ejemplos} de ubicacion, \textbf{dos} de medida y \textbf{dos} de dato sobre Almeria.}{7}

\BloomExercise{aplicar}{2}{Corrige frases}{Corrige estas frases para que sean matematicamente mas utiles: \emph{La playa es larga}; \emph{La Alcazaba esta lejos}; \emph{Hace calor en verano}.}{7}

        \begin{evidenciabox}
        \begin{itemize}
    \item Clasificar informacion urbana en ubicacion, medida y dato.
    \item Redactar ejemplos correctos con unidad cuando sea necesario.
    \item Elegir una forma simple de representar la informacion.
\end{itemize}
        \end{evidenciabox}

        \end{document}
